% !TEX root = ../../tir_monograph_v12.tex
\chapter{Complex Projective Geometry, Berry Holonomy, and Spin}
\label{ch:v12_complex_carrier}

\section{Input from the half-seam geometry}

Chapter~\ref{ch:v12_half_ln2} supplies three upstream objects:
\begin{equation}
\boxed{
\{N,S\},
\qquad
U(1)\cong S^1,
\qquad
\Sigma S^1\cong S^2.
}
\label{eq:v12_ch3_input}
\end{equation}
The present chapter realizes this sphere as the standard complex projective two-state geometry. No $SO(3)$ action is used to create the sphere.

\section{Complex two-state carrier}

Choose the two polar roles as a basis,
\begin{equation}
|N\rangle,\qquad |S\rangle.
\end{equation}
A normalized complex two-state vector has the form
\begin{equation}
|\psi\rangle
=
\alpha|N\rangle+\beta|S\rangle,
\qquad
|\alpha|^2+|\beta|^2=1.
\end{equation}
After removal of global phase,
\begin{equation}
|\psi\rangle\sim e^{i\lambda}|\psi\rangle,
\end{equation}
the ray space is
\begin{equation}
\boxed{
P(\mathbb C^2)
=
\mathbb{CP}^1
\cong
S^2.
}
\label{eq:v12_cp1_s2}
\end{equation}

A normalized representative may be written
\begin{equation}
\boxed{
|\psi(\vartheta,\varphi)\rangle
=
\cos\frac{\vartheta}{2}|N\rangle
+
e^{i\varphi}\sin\frac{\vartheta}{2}|S\rangle.
}
\label{eq:v12_spinor_spherical}
\end{equation}
The associated probability coordinate is
\begin{equation}
u=\sin^2\frac{\vartheta}{2},
\qquad
1-u=\cos^2\frac{\vartheta}{2}.
\end{equation}
Thus $u=1/2$ is exactly the equator $\vartheta=\pi/2$.

\section{Fubini--Study metric}

The standard Fubini--Study line element is
\begin{equation}
ds_{\rm FS}^2
=
\langle d\psi|d\psi\rangle
-
|\langle\psi|d\psi\rangle|^2.
\end{equation}
For Eq.~\eqref{eq:v12_spinor_spherical},
\begin{equation}
\boxed{
ds_{\rm FS}^2
=
\frac14
\left(
d\vartheta^2+\sin^2\vartheta\,d\varphi^2
\right).
}
\label{eq:v12_fs_qubit}
\end{equation}
Therefore the projective two-state space carries the round metric with Fubini--Study radius $1/2$:
\begin{equation}
\boxed{
\operatorname{Area}_{FS}(\mathbb{CP}^1)=\pi.
}
\end{equation}

\section{Berry connection and curvature}

In a standard gauge,
\begin{equation}
\mathcal A_B
=
i\langle\psi|d\psi\rangle
=
\frac{1-\cos\vartheta}{2}\,d\varphi
\end{equation}
up to orientation/sign convention. Hence
\begin{equation}
\boxed{
F_B=d\mathcal A_B
=
\frac12
\sin\vartheta\,d\vartheta\wedge d\varphi.
}
\label{eq:v12_berry_curvature}
\end{equation}
The integrated curvature is
\begin{equation}
\int_{S^2}F_B=2\pi,
\end{equation}
so
\begin{equation}
\boxed{
c_1
=
\frac1{2\pi}\int_{S^2}F_B
=
1.
}
\label{eq:v12_chern_one}
\end{equation}
This is the standard topological unit of the Berry phase on $\mathbb{CP}^1$.

\section{Balanced equatorial Berry holonomy}

For a latitude at fixed $\vartheta$, the Berry phase is one half of the enclosed solid angle, with sign fixed by orientation convention. In the equivalent population coordinate $\sigma$, the repository's information-spinor surface writes
\begin{equation}
\gamma_B(\sigma)
=
-2\pi(1-\sigma).
\end{equation}
At the balanced point,
\begin{equation}
\sigma=\frac12,
\end{equation}
one obtains
\begin{equation}
\boxed{
\gamma_B=-\pi\pmod{2\pi},
\qquad
e^{i\gamma_B}=-1.
}
\label{eq:v12_balanced_berry_sign}
\end{equation}

Thus the same half coordinate that maximizes the binary information measure sits on the nontrivial projective-sign Berry orbit.

\section{Euler--Berry selection of the primitive spin sector}

The archived formal module
\begin{center}
\texttt{archive/v7.9/full/22\_euler\_identity\_berry\_phase\_spin\_constraint\_v2\_4/}
\end{center}
records the following gate with status \texttt{FORMAL\_SYMBOLIC\_PASS}.

For spin sector $s$, the Berry phase of the primitive hemisphere loop is
\begin{equation}
\gamma_B=-s\,\Omega,
\qquad
\Omega=2\pi,
\end{equation}
up to orientation convention. Requiring the nontrivial Euler projective sign,
\begin{equation}
e^{i\gamma_B}=-1,
\end{equation}
gives
\begin{equation}
2\pi s=\pi\pmod{2\pi}.
\end{equation}
The minimal positive solution is
\begin{equation}
\boxed{s=\frac12.}
\label{eq:v12_spin_half_selection}
\end{equation}

The doubled loop closes the lifted spinor:
\begin{equation}
\boxed{
2\pi\mapsto-I,
\qquad
4\pi\mapsto+I.
}
\label{eq:v12_spin_double_cover}
\end{equation}

This is a formal/projective spin-selection result. It is not promoted here to a claim that all physical spin assignments follow from this chapter alone.

\section{Polar closure}

For the spin-$1/2$ realization, a distinction generator may be chosen as $\sigma_z$:
\begin{equation}
\sigma_z|N\rangle=+|N\rangle,
\qquad
\sigma_z|S\rangle=-|S\rangle.
\end{equation}
Hence the spinorial geometry reproduces the same polar pair that entered upstream as the two roles of the primitive relation:
\begin{equation}
\boxed{
\{N,S\}_{R}
\longleftrightarrow
\{N,S\}_{\rm spin}.
}
\label{eq:v12_polar_closure}
\end{equation}
This is a closure certificate, not a new parent edge; it therefore does not create a cycle in the dependency DAG.

\section{Density-matrix and Bloch representation}

For a normalized pure ray,
\begin{equation}
\rho_\psi=|\psi\rangle\langle\psi|
=
\frac12
\left(
I+\mathbf r\cdot\boldsymbol\sigma
\right),
\qquad
|\mathbf r|=1,
\end{equation}
with
\begin{equation}
\mathbf r
=
(
\sin\vartheta\cos\varphi,\,
\sin\vartheta\sin\varphi,\,
\cos\vartheta
).
\end{equation}
For a general normalized qubit state,
\begin{equation}
|\mathbf r|\le1,
\end{equation}
and the trace-one Hermitian affine hull is
\begin{equation}
\boxed{
\mathcal A_2
=
\frac12I+\operatorname{Herm}_0(2).
}
\label{eq:v12_affine_qubit_hull}
\end{equation}

\section{Operator coordinates and the three-real-dimensional carrier}

Every Hermitian $2\times2$ operator has the unique Pauli decomposition
\begin{equation}
A
=
a_0I
+
a_x\sigma_x
+
a_y\sigma_y
+
a_z\sigma_z,
\qquad
a_\mu\in\mathbb R.
\end{equation}
The distinction-generating part is
\begin{equation}
\boxed{
\operatorname{Herm}_0(2)
=
\operatorname{span}_{\mathbb R}
\{\sigma_x,\sigma_y,\sigma_z\}
\cong
\mathbb R^3.
}
\label{eq:v12_herm0_preview}
\end{equation}

The Hilbert--Schmidt form
\begin{equation}
\langle A,B\rangle_{\rm HS}
=
\frac12\operatorname{Tr}(AB)
\end{equation}
gives
\begin{equation}
\frac12\operatorname{Tr}(\sigma_i\sigma_j)
=
\delta_{ij}.
\end{equation}
Chapter~\ref{ch:v12_c2_herm0_r3} develops the resulting affine relation carrier.

\section{Rotation covariance comes last}

For $U\in SU(2)$,
\begin{equation}
A\mapsto UAU^\dagger
\end{equation}
preserves trace and the Hilbert--Schmidt metric. On coefficient vectors,
\begin{equation}
\boxed{
\operatorname{Ad}:SU(2)\to SO(3),
\qquad
SU(2)/\{\pm I\}\cong SO(3).
}
\label{eq:v12_su2_so3}
\end{equation}

The dependency order is therefore explicit:
\begin{equation}
\boxed{
S^1+\{N,S\}
\to
S^2
\cong
\mathbb{CP}^1
\to
(g_{\rm FS},F_B)
\to
s=\frac12
\to
SU(2)/\{\pm I\}
\cong
SO(3).
}
\label{eq:v12_ch3_dependency}
\end{equation}
In particular, $SO(3)$ is downstream of the sphere and is never used to manufacture it.

\section{Chapter output}

The chapter exports
\begin{equation}
\boxed{
\mathbb C^2
\to
\mathbb{CP}^1\cong S^2
\to
(g_{\rm FS},F_B,c_1=1)
\to
s=\frac12
\to
\operatorname{Herm}_0(2)\cong\mathbb R^3.
}
\end{equation}

The canonical source-order owner is
\begin{center}
\texttt{TIR/foundations/TIR\_CANONICAL\_DERIVATION\_SPINE\_V0\_1.md}.
\end{center}
The archived Euler--Berry theorem is cited as a promoted formal proof source with its original epistemic status preserved.
